\appendixsection{实验复现与结果文件说明}\label{sec:appendix-reproducibility}

本附录列出本文新增实验与论文图表的主要复现入口，便于后续检查、答辩说明和结果追溯。所有新增对比实验均以结构化JSON文件作为结果来源，论文中的算法对比表格和可视化图均从这些结果文件生成。

\appendixsubsection{主要实验脚本}

\begin{itemize}
  \item \texttt{experiments/compare\_algorithms.py}：在相同场景、相同随机种子和相同评价预算下运行本文MOPSO-DT、原始MOPSO-DT、NSGA-II-DT、MOEA/D-DT、SPEA2-DT和随机搜索。
  \item \texttt{experiments/boundary\_analysis.py}：在L形复杂区域中生成边界任务点，并计算Overall ECR、Boundary ECR、Boundary Gap、HV和运行时间等指标。
  \item \texttt{tools/generate\_comparison\_figures.py}：读取结构化结果文件并生成算法前沿叠加图、指标柱状图、质量-效率气泡图、边界覆盖图和拐点部署对比图。
\end{itemize}

\appendixsubsection{结构化结果文件}

\begin{itemize}
  \item \texttt{results/algorithm\_comparison.json}：保存Challenging场景下多算法三随机种子的完整非支配档案、评价指标和运行时间。
  \item \texttt{results/boundary\_analysis.json}：保存边界专项实验中各方法的边界覆盖指标、代表性部署方案和运行统计。
  \item \texttt{TongjiThesis-1.4.0/figures/}：保存论文当前编译使用的图表文件。
\end{itemize}

\appendixsubsection{正式运行命令}

\begin{quote}
\small\ttfamily\raggedright
python experiments\textbackslash compare\_algorithms.py\\
\hspace*{1em}--scenario challenging --seeds 2026 2027 2028\\[2pt]
python experiments\textbackslash boundary\_analysis.py\\
\hspace*{1em}--seeds 2026 2027 2028\\[2pt]
python tools\textbackslash generate\_comparison\_figures.py
\end{quote}
